\subsection*{Inherited results and the present additions}
The complete localized-clearance continuation, including the source
renewal, deep-clearance payment, all-label variation, conservative
return tower and endpoint-tested weak operators, is retained in
Sections~\ref{sec:lc-length}--\ref{sec:wo-endpoint-tests} at its actual
strength. In particular its deep-clearance theorem leaves a shallow
source with the same full nongrazing pool. No independent recurrence
or hidden collision deletion is introduced here.

Our first new estimate addresses a selected clearance flight without
discarding the two endpoint layers. If $D_{m,R;j,\delta}(a,v)$ is the
complete physical inverse-derivative sum with initial and terminal
momenta $a,v$ and a clearance witness at flight $j$, then
Theorem~\ref{thm:ce-kernel} proves
\[
 D_{m,R;j,\delta}(a,v)
 \le C_{\rm end}\{\delta+\theta_{\rm r}^{\max(j,m-1-j)}\}.
\]
The first flight is paid by its exact weighted image length.
Time reversal compares two sums on the same physical trajectories;
there is no product-independence claim. The resulting critical
sector estimate, Theorem~\ref{thm:ce-critical}, is summed over all
actual words and witnesses. It controls the unscaled critical
coefficient, with an explicit exponentially decaying remainder,
and does not supply the missing collision factor $m^{-2}$.

Our second new result inverts the exact roof-interval event. For a
union $I$ of at most $J$ finite intervals, the full marked transform
has an absolutely convergent inverse after multiplication by
$\widehat{\one_I}(-s)$. Its cutoff at $|s|\le S$ has error at most
$2JD/(\pi S)$, where $D$ is the proved total variation budget.
Theorem~\ref{thm:ui-inversion} and Corollary~\ref{cor:ui-all-labels}
retain the entire lattice and occupation torus and allow the sum
over all exact labels without a label-count multiplier.
Proposition~\ref{prop:ui-conditional} uses the actual microscopic
return probability; an absolute error is never divided by an
unproved typical probability. Proposition~\ref{prop:ui-raw-bridge}
states a sufficient ordered raw-height criterion. Its fixed-window
terms are evaluated using the retained native 87 theorem in
Proposition~\ref{prop:ui-inherited-bridge}; the normalized local
derivative charge remains a substantive sufficient condition.

Finally Section~\ref{sec:sm-section} supplies explicit constants
for the original moving section and checks the precise strong
multiplier hypotheses. The mechanism is inherited; the new
quantitative supplement bounds the defining root derivative,
the non-grazing section cap and the coordinate displacement.
The full four-record spectral results for fixed bands already
exist in the historical native proof chain. Bounds uniform on
growing bands remain a different assertion, but they are not
required for the stated scalar implication through the retained
fixed-window theorem. The exact arithmetic reference retains all
the previously specified peripheral and near-peripheral branches.
